% 1. ročník MAT – příprava 102; inspirace Realisticky.cz 020501–020506, úlohy nově sestavené.
\documentclass[11pt,a4paper]{article}
\input{../../sablony/pripravy/priprava.sty}
\usepackage{tikz}
\usepackage{pgfplots}
\pgfplotsset{compat=1.18}
\usetikzlibrary{arrows.meta}
\nastavPripravu{1. ročník}{Kvadratické funkce}{102}{Souhrnné procvičování kvadratických funkcí}
\begin{document}
\blok{Souhrnné procvičování – ne nová teorie}
Po hodinách 96–101 už víme, jak poznat kvadratickou funkci, zakreslit $ax^2$, posunout parabolu, doplnit na čtverec, najít vrchol, obor hodnot a jednoduché průsečíky s osami. Dnes studenti vybírají úlohy ze \textbf{samostatného pracovního listu 102}, který je rozdělený na A (základy), B (konstrukce) a C (výzva).
\blok{Organizace 45 minut}
\begin{enumerate}[label=\arabic*.]
\item 5 min: každý žák označí dvě oblasti, které mu dělají potíže.
\item 7 min: učitel ukáže jeden náročnější krok doplnění na čtverec nebo odečtení vrcholu.
\item 25 min: samostatná práce na pracovním listu; doporučené minimum \emph{jedna úloha z A, tři z B a jedna s grafem}.
\item 8 min: společná kontrola postupu a diskuse, proč různé tvary mohou popisovat stejnou funkci.
\end{enumerate}
\blok{Vzorová kontrola 1: z předpisu ke grafu}

\begin{prikladReseniBox}
Upravte $f(x)=-x^2+4x-3$ na vrcholový tvar. Určete všechny informace potřebné pro rychlý náčrt.
\tcblower
\[
\begin{aligned}
f(x)&=-(x^2-4x)-3 \\
&=-(x-2)^2+1 \\
V&=[2;1],\quad H(f)=(-\infty;1\rangle \\
f(0)&=-3\quad\Rightarrow\ [0;-3] \\
f(x)=0&\iff(x-1)(x-3)=0 \\
P_{x,1}&=[1;0],\quad P_{x,2}=[3;0]
\end{aligned}
\]

\end{prikladReseniBox}
\blok{Vzorová kontrola 2: najdi chybu}

\begin{zadaniReseniBox}
Žák tvrdí: „Funkce $g(x)=2(x+3)^2-4$ má vrchol $[3;-4]$.“ Opravte jeho úvahu a vysvětlete znaménko uvnitř závorky.
\tcblower
\[
\begin{aligned}
(x+3)^2=0&\iff x=-3 \\
V_g&=[-3;-4] \\
H(g)&=\langle-4;+\infty)
\end{aligned}
\]
Kritický argument je $x=-3$; při $x=3$ je výraz $(x+3)^2$ roven $36$, nikoliv nule.
\end{zadaniReseniBox}
\blok{Samostatná práce se sbírkou}

Žákovský pracovní list \texttt{102\_PL\_KvadratickeFunkce.pdf} obsahuje \textbf{18 úloh} postupně od určování vlastností po rekonstrukci předpisu a vzájemné průsečíky grafů. Pro úlohy s náčrtem jsou připojené soustavy souřadnic. \par
Úplný učitelský klíč je v samostatném souboru \texttt{102\_PL\_KvadratickeFunkce\_reseni.pdf}.
\blok{Volitelná rezerva}
Když žáci dokončí povinné minimum, mohou řešit zadání C z listu. Výslovně přenech volbu postupů studentům: vrcholový tvar, rozklad i dosazování kontrolních bodů jsou rovnocenné vhodně použité strategie.
\blok{Domácí úkol}
\begin{zadaniReseniBox}[title={Domácí příprava}]
\begin{enumerate}[label=\alph*),itemsep=3pt]
\item Pro $f(x)=2x^2$ určete $f(-2)$, vrchol a obor hodnot.
\item U $g(x)=-(x-2)^2+5$ určete vrchol, osu a maximum.
\item Upravte $x^2+8x+12$ na vrcholový tvar.
\item U $h(x)=x^2-2x-8$ určete nulové body.
\item Vypočtěte vrchol a $H(f)$ pro $f(x)=-2x^2-4x+1$.
\item \textbf{Výzva:} Určete $a$, má-li $f(x)=a(x-1)^2-2$ procházet bodem $[3;6]$.
\end{enumerate}
\tcblower
\begin{enumerate}[label=\alph*),itemsep=4pt]
\item
\[
\begin{aligned}
f(-2)&=8,\quad V=[0;0] \\
H(f)&=\langle0;+\infty)
\end{aligned}
\]

\item
\[
\begin{aligned}
V&=[2;5],\quad\text{osa }x=2,\quad\max g=5
\end{aligned}
\]

\item
\[
\begin{aligned}
x^2+8x+12&=(x+4)^2-4
\end{aligned}
\]

\item
\[
\begin{aligned}
x^2-2x-8&=(x-4)(x+2) \\
K&=\{-2;4\}
\end{aligned}
\]

\item
\[
\begin{aligned}
-2x^2-4x+1&=-2(x+1)^2+3 \\
V&=[-1;3],\quad H(f)=(-\infty;3\rangle
\end{aligned}
\]

\item
\[
\begin{aligned}
6&=a(3-1)^2-2 \\
8&=4a\ \Rightarrow\ a=2
\end{aligned}
\]

\end{enumerate}
\end{zadaniReseniBox}
\poznamka{Nechávej žáky vysvětlovat \emph{proč} volili postup. Před kvadratickými rovnicemi je důležitější správný vrchol a graf než rychlost výpočtu. Neměň úlohy A/B/C na povinnou písemku všech 18 příkladů.}
\end{document}
